\section{Notation and Learning Setting}

This chapter fixes the technical vocabulary used by the later methodology and evaluation chapters. The goal is not to survey every defense or attack, but to define the optimization setting in which adversarial examples, multi-attack training, group reweighting, clustering, and multi-seed evaluation are used. The notation is kept deliberately close to the notation used in Chapters~\ref{chap:methodology} and~\ref{chap:experimental_setup}.

We consider supervised classification with input space $\mathcal{X}\subseteq\mathbb{R}^{d}$ and label space $\mathcal{Y}=\{1,\ldots,C\}$. A training set is denoted by
\[
\mathcal{D}=\{(x_i,y_i)\}_{i=1}^{n},
\]
where each pair is sampled from an underlying data distribution $P_{XY}$. A classifier with parameters $\theta$ maps an input to logits
\[
f_\theta:\mathcal{X}\rightarrow\mathbb{R}^{C},
\]
and predicts $\hat{y}=\arg\max_{c} f_\theta(x)_c$. The softmax transformation converts logits into class probabilities,
\[
p_\theta(y=c\mid x)=
\frac{\exp(f_\theta(x)_c)}
{\sum_{j=1}^{C}\exp(f_\theta(x)_j)}.
\]

\subsection{Deep Neural Networks and Supervised Classification}

Deep neural networks are used here as parameterized classifiers rather than as objects of architectural study. Their relevance to this report is that the same differentiable structure used for gradient-based learning also exposes gradients with respect to the input. These input gradients are the basic computational object behind many first-order adversarial attacks, including FGSM and PGD \cite{goodfellow2015explaining,madry2019towards}.

For a labeled example $(x,y)$, the standard classification loss is cross-entropy,
\[
\ell(f_\theta(x),y)=-\log p_\theta(y\mid x).
\]
Empirical risk minimization (ERM) fits the model by minimizing average training loss:
\[
\min_{\theta}
\frac{1}{n}\sum_{i=1}^{n}\ell(f_\theta(x_i),y_i).
\]
ERM is the clean-data reference point for the rest of the report. It optimizes performance on sampled training examples, but it does not require the prediction to remain stable inside a local neighborhood of each input. Adversarial robustness adds precisely this local stability requirement.

\subsection{Loss Functions and Gradient-Based Optimization}

In gradient-based training, the model parameters are updated using derivatives of the loss with respect to $\theta$. In gradient-based attacks, the adversary instead uses derivatives with respect to $x$. A small perturbation $\delta$ can therefore be selected to increase the loss while remaining visually or semantically close to the original input:
\[
\arg\max_c f_\theta(x)_c
\neq
\arg\max_c f_\theta(x+\delta)_c,
\quad
\|\delta\|_p\leq \varepsilon .
\]
This mismatch between small input displacement and large prediction change was established by early adversarial-example work and remains the core failure mode studied in adversarial training \cite{szegedy2014intriguing,goodfellow2015explaining}.

The later chapters use this notation to distinguish attack algorithms from training methods. For example, PGD-$\ell_\infty$ is an attack used to generate perturbed examples, whereas PGD-AT is a training method that repeatedly trains on PGD-generated examples. Likewise, DDN-$\ell_2$ is an attack, while DDN-AT is a single-source adversarial training baseline. This distinction matters because Chapter~\ref{chap:methodology} compares training strategies while holding the source attack set fixed.

\section{Adversarial Examples and Threat Models}

An adversarial example is an input deliberately perturbed to induce an incorrect model prediction while preserving the ground-truth semantic label. Given a clean example $(x,y)$, an adversarial example has the form
\[
x^{\mathrm{adv}}=x+\delta .
\]
In an untargeted attack, the adversary searches for a perturbation that makes the classifier predict any class other than $y$. In a targeted attack, the adversary instead tries to force prediction of a chosen target class $y_{\mathrm{target}}\neq y$. This report primarily uses untargeted robustness evaluation, with TPGD treated according to the TRADES-style implementation described in Chapter~\ref{chap:experimental_setup}.

\subsection{Norm-Bounded Perturbation Sets}

Most image-classification robustness evaluations constrain $\delta$ by an $\ell_p$ budget. The feasible perturbation set is
\[
\mathcal{S}_p(\varepsilon)=
\{\delta\in\mathbb{R}^{d}:\|\delta\|_p\leq \varepsilon\},
\]
with $p$ usually chosen as $2$ or $\infty$. The perturbed input is also clipped to the valid input range, so the effective search region is
\[
x^{\mathrm{adv}}\in (x+\mathcal{S}_p(\varepsilon))\cap[0,1]^d .
\]
The label-preservation assumption states that the semantic label should not change inside this small perturbation set. Under that assumption, a misclassification at $x^{\mathrm{adv}}$ is treated as a robustness failure rather than a change in the underlying task.

For a fixed model and loss, adversarial example generation is often written as an inner maximization:
\[
\delta^\ast \in
\arg\max_{\delta\in\mathcal{S}_p(\varepsilon)}
\ell(f_\theta(x+\delta),y).
\]
Different attacks approximate this maximization in different ways. Some follow the input gradient, some optimize a margin-based objective, and others estimate a nearby decision boundary. These algorithmic differences are important because robustness against one perturbation family or loss surrogate does not necessarily transfer to another \cite{tramer2019adversarial,maini2020adversarial,croce2020reliable}.

For attack condition $e$, let $\mathcal{A}_e$ denote the corresponding adversarial example generator. The generated input can be written as
\[
x^{\mathrm{adv}}_e=\mathcal{A}_e(x,y,f_\theta),
\]
and the associated attack-domain risk is
\[
R_e(\theta)=
\mathbb{E}_{(x,y)\sim P_{XY}}
\left[
\ell(f_\theta(x^{\mathrm{adv}}_e),y)
\right].
\]
This notation supports the later multi-attack setting: a method can reduce risk for one source attack while leaving another attack-domain risk high. The training methods in this report therefore compare not only average adversarial risk, but also how training emphasis is allocated across attack-induced conditions and latent groups.

\subsection{Whitebox, Graybox, and Blackbox Access}

The perturbation set specifies what changes are allowed in input space, while the threat model specifies what information the adversary can use. In the whitebox setting, the attacker has access to the model architecture, parameters, loss, and gradients. First-order attacks such as FGSM and PGD are whitebox attacks because they use $\nabla_x \ell(f_\theta(x),y)$ directly. Whitebox evaluation is the standard setting for testing gradient-based robustness and is also the natural setting for adversarial training, where the current model generates its own adversarial examples \cite{madry2019towards,athalye2018obfuscated}.

In the graybox setting, the adversary has partial information. A common case is transfer evaluation: adversarial examples are generated on a surrogate model $f_{\theta_s}$ and evaluated on a target model $f_{\theta_t}$. This tests whether perturbations built from one loss landscape remain harmful on another. Graybox transfer is relevant to the attacks-as-domains view because it changes the mechanism that generates the adversarial distribution without necessarily changing the data distribution itself.

In the blackbox setting, the adversary does not access model parameters or gradients and instead relies on queries. Score-based attacks use probabilities or logits, while decision-based attacks use only predicted labels. Blackbox attacks are useful for detecting gradient masking, where a defense appears strong because gradients are uninformative rather than because the classifier has learned genuinely robust decision boundaries \cite{athalye2018obfuscated,andriushchenko2019square}. The main experiments in this report focus on the specified whitebox evaluation suite and separate stress checks, while the blackbox setting remains part of the broader threat-model vocabulary.

These definitions prepare the move from individual attacks to training objectives. Once attacks are viewed as generators of distinct adversarial risks, multi-attack and group-aware training can be framed as choices about which risks are optimized, how they are weighted, and how hidden failure modes are discovered.
